% tonks_paper04_v1.tex
% Generalised hydrodynamics on Riemannian rings with quenched
% binary mass disorder.
% v1 (2026-05-13): initial skeleton
% v1.5 (2026-05-15): erratum to Eq. (10), (12), and linearisation:
%   the dressed-velocity numerator must use η·ū (not σ·ū), since σ
%   has units of length and ū of velocity. The dimensionally correct
%   form is v_eff = (v - η ū)/(1-η), equivalent to (v - σ p̄)/(1-σρ)
%   where p̄ = ρū is the momentum density. The analytical c_s in
%   Eq. (14) is unaffected. Verified against discrete-velocity solver
%   (Paper V): with corrected dressing, solver matches analytic c_s
%   to 2% across η ∈ [0.01, 0.20] and m_max ∈ [2, 10].

\pdfminorversion=7
\documentclass[aps,pre,twocolumn,superscriptaddress,longbibliography,floatfix]{revtex4-2}

\usepackage{amsmath,amssymb,amsfonts,amsthm}
\usepackage{graphicx}
\usepackage{bm}
\usepackage{hyperref}
\usepackage{physics}
\usepackage{mathtools}

\begin{document}

\title{Generalised hydrodynamics on Riemannian rings with quenched
binary mass disorder}

\author{Hector Medel}
\email{hjmedel@gmail.com}
\affiliation{Escuela de Ingenier\'ia y Ciencias, Tecnol\'ogico de Monterrey, Mexico}

\date{\today}

\begin{abstract}
We construct the generalised hydrodynamic (GHD) equation for
the hard-rod gas on a closed one-dimensional Riemannian
curve $(M,g)$ and extend it to a binary mass-disordered mixture
with quenched disorder. The geometric content of the curved
problem is shown to reduce, by the arc-length theorem of
companion paper~\cite{PaperI}, to that of the flat Tonks gas of
perimeter $L=\oint\sqrt{g}\,d\varphi$, with the metric entering
only through observable averages on the manifold. For the
binary mixture $m_i\in\{m_L,m_H\}$, we construct a two-species
GHD-Boltzmann hybrid equation in which the inter-species
elastic collision rate sets a thermalisation kernel of strength
$\tau_{\rm break}^{-1}$. Linearised about equilibrium, the equation predicts a
sound-mode dispersion that matches a dedicated event-driven MD
sweep (extending the $m_{\max}=5$ result of companion
paper~\cite{PaperII} to $m_{\max}\in\{5,10,100\}$) within
$\sim 4\%$ across a 20-fold range in mass disparity. At finite
amplitude with step velocity initial conditions, the linearised
analysis predicts a density-mode ratio
$|\hat\rho_n|/|\hat\rho_1|=1/n$ for odd $n$, $m_{\max}$-%
independent, recovered to $5\%$ in the data of companion
paper~\cite{PaperIII}; this provides the mechanism for the
empirical density-mode universality reported there. The
even-mode amplitudes and their slow drift across $m_{\max}$
are not derived from first principles here and are left as
open work.
\end{abstract}

\maketitle

\section{Introduction}\label{sec:intro}

Generalised hydrodynamics (GHD) provides an exact Euler-scale
description of the hydrodynamic limit of integrable many-body
systems in one dimension
\cite{DoyonSpohn2017,Doyon2017Lectures,BertiniGhdReview2025,
DenardisDoyon2018,Bulchandani2018,GHDperspective2025}. For the
classical hard-rod gas, it was first written by Boldrighini,
Dobrushin and Sukhov \cite{BDS1983} and recently revisited
in connection with diffusive corrections \cite{HRDiffusive2023}
and absence of shocks \cite{HubnerDoyon2024}. Two extensions
of the standard framework are relevant for the present work:

\emph{(i)} The geometric setting. GHD has been formulated for
gases on the real line and on a periodic ring with translation
invariance; the case of a non-trivially curved one-dimensional
manifold $(M,g)$ has not, to our knowledge, been treated
explicitly. The closest result is the ``geometric viewpoint''
of Doyon, Spohn and Yoshimura \cite{DoyonSpohnYoshimura2017},
which introduces a state-dependent metric on the rapidity
space rather than on physical space.

\emph{(ii)} Integrability breaking. A growing body of work on
weak-integrability-breaking GHD \cite{DurninDoyon2021,
WeakBreaking2023, BastianelloBoltzmann2020} treats stochastic
or annealed perturbations of integrable models, deriving
Boltzmann-type collision integrals that quantify
thermalisation. The case of \emph{quenched} disorder ---
relevant for a binary mass mixture where each realisation
carries a fixed random assignment of masses --- has not been
systematically addressed.

In this work we combine the two extensions. Section
\ref{sec:setup} recalls the setup of the hard-rod gas on $(M,g)$
and the arc-length reduction theorem of
Ref.~\cite{PaperI}. Section \ref{sec:ghdmg} writes the GHD
equation in covariant form on $(M,g)$ and proves its equivalence
to the flat-space equation on the ring of length $L$.
Section \ref{sec:twospecies} constructs a two-species GHD
equation for binary mass disorder, postulating a standard
Boltzmann form for the inter-species kernel and motivating
the structure of the dressed velocity. Section \ref{sec:linear} compares the linearised
equation against the sound-wave data of
Ref.~\cite{PaperII}. Section \ref{sec:nonlinear} addresses the
nonlinear regime and the $m_{\max}$-universality observed in
Ref.~\cite{PaperIII}. Section \ref{sec:discussion} situates
the result and lists open work.

The main contributions are:
\emph{(i)} a covariant GHD equation on Riemannian 1D rings
and a proof that physical-space curvature enters only through
the perimeter $L$ at the Euler scale;
\emph{(ii)} a postulated two-species GHD-Boltzmann equation for
the binary-mass hard-rod gas with quenched disorder, with a
standard Boltzmann-form inter-species kernel of strength
$\tau_{\rm break}^{-1}$ (the explicit derivation from
microscopic crossing dynamics is left to future work);
\emph{(iii)} quantitative consistency with the linear-regime
sound-wave restoration of Ref.~\cite{PaperII};
\emph{(iv)} a mechanism for the $m_{\max}$-universality of the
density-mode spectrum at finite amplitude reported in
Ref.~\cite{PaperIII}.

\section{Setup and arc-length reduction}\label{sec:setup}

We consider $N$ hard rods of diameter $\sigma$ on a smooth
closed curve $\gamma:S^1\to\mathbb{R}^d$ with arc-length element
$ds=\sqrt{g(\varphi)}\,d\varphi$ and perimeter
$L=\oint\sqrt{g}\,d\varphi$. Half the rods have mass $m_L=1$ and
half have mass $m_H = m_{\max}>1$; the assignment is fixed per
realisation and is quenched. The single-realisation Hamiltonian
is
\begin{equation}
H = \sum_{i=1}^N \frac{p_i^2}{2 m_i}
  + \sum_{i<j} V_{\rm rod}(s_i-s_j),
\label{eq:H}
\end{equation}
with $V_{\rm rod}$ the hard-rod pair potential of range $\sigma$.

By the arc-length theorem of Ref.~\cite{PaperI}, the canonical
configurational integral on $(M,g)$ in arc-length coordinate is
identical to that of the flat Tonks gas on a circle of
perimeter $L$. Dynamically, the equations of motion in
arc-length coordinate are those of the flat gas as well: each
rod moves at constant velocity between point elastic collisions,
and the collision rule is mass-dependent,
\begin{equation}
v_i' = \frac{m_i - m_j}{m_i + m_j}\,v_i
     + \frac{2 m_j}{m_i + m_j}\,v_j.
\label{eq:elcoll}
\end{equation}
The geometry enters only through (a) the perimeter $L$ and
(b) the relation between arc-length and intrinsic coordinate
observables; the bulk dynamics is geometrically trivial.

\section{GHD on $(M,g)$}\label{sec:ghdmg}

For the integrable equal-mass case ($m_L=m_H$), the Euler-scale
GHD equation in arc-length coordinate $s$ reads
\begin{equation}
\partial_t \rho_p(s,v,t)
+ \partial_s\!\left[v^{\rm eff}(s,v,t)\,\rho_p(s,v,t)\right] = 0,
\label{eq:ghd}
\end{equation}
where $\rho_p(s,v,t)$ is the rapidity (asymptotic velocity)
density and $v^{\rm eff}$ is the dressed velocity. For the
hard-rod gas of total density $\rho(s,t)=\int dv\,\rho_p(s,v,t)$,
the standard result is \cite{DoyonSpohn2017,BDS1983}
\begin{equation}
v^{\rm eff}(s,v,t) = \frac{v - \eta(s,t)\,\bar u(s,t)}{1 - \eta(s,t)},
\quad
\bar u(s,t) = \rho^{-1}\!\int\!dv\,v\,\rho_p,
\label{eq:veff}
\end{equation}
with $\bar u(s,t)$ the local number-weighted mean velocity and
$\eta = \sigma\rho$ the local packing fraction (dimensionless).
The combination $\eta\bar u$ is a velocity, so each term in the
numerator has consistent units;
equivalently, $\eta\bar u = \sigma\bar p$ where
$\bar p = \int\!dv\,v\,\rho_p = \rho\bar u$ is the momentum density
per unit length of the rapidity distribution.

\subsection{Covariant form on $(M,g)$}

The arc-length theorem of Ref.~\cite{PaperI} states that the
hard-rod dynamics on $(M,g)$ is identical to the flat gas in the
arc-length coordinate. Consequently the GHD equation on the
manifold is given by Eq.~\eqref{eq:ghd} with $s$ the arc-length
coordinate and $L=\oint\sqrt{g}\,d\varphi$ the periodicity. In
the intrinsic coordinate $\varphi$ with arc-length element
$\sqrt{g(\varphi)}d\varphi$, the same equation reads
\begin{equation}
\partial_t \tilde\rho_p(\varphi,v,t)
+ \frac{1}{\sqrt{g}}\partial_\varphi\!\left[v^{\rm eff}\,\tilde\rho_p\right]
= 0,
\label{eq:ghdmg}
\end{equation}
with $\tilde\rho_p(\varphi,v,t) = \sqrt{g(\varphi)}\,\rho_p(s(\varphi),v,t)$
the density per unit $\varphi$. The factor $\sqrt{g}$ is the
Jacobian relating the two parametrisations.

A direct consequence is the
\emph{Euler-scale geometric invariance}: at the Euler scale, the
GHD dynamics on $(M,g)$ depends on the metric $g(\varphi)$ only
through its integral $L$. This statement applies to bulk
observables; the relation between intrinsic-coordinate and
arc-length-coordinate observables of course retains the
Jacobian dependence on $g(\varphi)$. We will recover this
statement empirically in
Sec.~\ref{sec:nonlinear} against the curved-geometry data of
Ref.~\cite{PaperIII}.

\section{Two-species GHD with binary mass disorder}
\label{sec:twospecies}

For the binary mixture the asymptotic velocity is no longer
conserved per particle: an elastic collision between unequal
masses redistributes velocity via Eq.~\eqref{eq:elcoll},
breaking the integrability of the equal-mass case. There is
therefore no rapidity in the strict GHD sense. We work instead
with the species-resolved velocity distribution
$\rho_\alpha(s,v,t)$ on phase space ($\alpha\in\{L,H\}$),
normalised so that $\int\!dv\,\rho_\alpha = \rho_\alpha^{\rm
tot}(s,t)$ is the local density of species $\alpha$. The
relation to the equal-mass GHD framework
(Sec.~\ref{sec:ghdmg}) is that $\rho_\alpha$ reduces to the
single-species rapidity density $\rho_p$ in the integrable
limit $m_L=m_H$.

\subsection{GHD-Boltzmann equation}

The species-resolved hydrodynamic equation has the form
\begin{equation}
\partial_t \rho_\alpha + \partial_s\!\left[v^{\rm eff}_\alpha\,\rho_\alpha\right]
= \mathcal{C}_\alpha[\rho_L,\rho_H](s,v,t),
\label{eq:ghdboltz}
\end{equation}
where $v^{\rm eff}_\alpha$ is the dressed velocity in the binary
mixture (to be determined below) and
$\mathcal{C}_\alpha[\rho_L,\rho_H]$ is a Boltzmann-style
collision operator capturing inter-species velocity exchange.

For the elastic collision Eq.~\eqref{eq:elcoll}, the standard
Boltzmann collision integral for a binary classical
hard-rod gas in one dimension is
\begin{multline}
\mathcal{C}_\alpha(s,v,t)
= \sigma\!\int\!dv'\,|v-v'|\\
\times \big[\rho_\alpha(s,v_*,t)\rho_\beta(s,v'_*,t)
           - \rho_\alpha(s,v,t)\rho_\beta(s,v',t)\big],
\label{eq:Calpha}
\end{multline}
where $\beta\neq\alpha$, $\sigma$ is the rod diameter setting
the 1D ``cross section'', and $(v_*,v'_*)$ is the pre-collision
pair that yields $(v,v')$ after the elastic update
Eq.~\eqref{eq:elcoll}. The Jacobian of the velocity-space
transformation $(v,v')\!\to\!(v_*,v'_*)$ equals unity because
elastic two-body collisions preserve phase-space volume.
Only inter-species collisions ($\alpha\neq\beta$) appear in
Eq.~\eqref{eq:Calpha}: equal-mass intra-species collisions
(L--L or H--H) reduce to a velocity swap and so preserve the
species-resolved velocity multiset, contributing to the dressed
streaming Eq.~\eqref{eq:veffalpha} but not to thermalisation.
The factorisation of pair distributions on the right-hand
side is the molecular-chaos closure (Stosszahlansatz), which
breaks microscopic time-reversal symmetry at the coarse-grained
level and is the source of the H-theorem discussed below.

\subsection{Dressed velocity in the binary mixture: postulate}

For the dressed velocity in the binary mixture we postulate
\begin{equation}
v^{\rm eff}_\alpha(s,v,t)
= \frac{v - \eta(s,t)\,\bar u(s,t)}{1 - \eta(s,t)},
\label{eq:veffalpha}
\end{equation}
with $\eta(s,t)=\sigma\rho^{\rm tot}(s,t)$ the local packing
fraction, $\rho^{\rm tot}=\rho_L^{\rm tot}+\rho_H^{\rm tot}$,
and $\bar u$ the number-weighted mean velocity
\begin{equation}
\bar u(s,t)
= \frac{\int dv\,v\,(\rho_L + \rho_H)}{\rho^{\rm tot}}.
\label{eq:ubar}
\end{equation}
The form Eq.~\eqref{eq:veffalpha} is dimensionally a velocity
in each term of the numerator (since $\eta$ is dimensionless
and $\bar u$ has velocity units), and reduces to
Eq.~\eqref{eq:veff} in the equal-mass limit.
The structural motivation is that the geometric scattering shift
on each crossing event is $\sigma$ regardless of the masses of
the two colliding rods: the modification of effective
velocity by collisions depends only on the rate of crossing
events, which is set by the relative number density and not by
mass weighting \cite{DoyonSpohn2017,BDS1983}. The mass-weighted
mean velocity does appear elsewhere, in the momentum-balance
equation of the Chapman--Enskog reduction
(Sec.~\ref{sec:linear}), but not in the dressed velocity. The mass dependence enters
entirely through the collision integral $\mathcal{C}_\alpha$.
We treat Eq.~\eqref{eq:veffalpha} as a working assumption and
test it via the quantitative predictions of
Sec.~\ref{sec:linear}: if the postulate is incorrect, the sound
speed will not reproduce the molecular-dynamics data. A
microscopic derivation of Eq.~\eqref{eq:veffalpha} from the
binary-mixture crossing statistics is left as follow-up work;
we note that the classical Jepsen mapping itself does not
straightforwardly extend to unequal masses, because the elastic
binary collision Eq.~\eqref{eq:elcoll} is not a velocity swap,
so the derivation will require a modified approach.

\subsection{Inter-species thermalisation rate}

For the binary fraction $\phi=1/2$ at packing $\eta$ and
temperature $k_BT$,
\begin{equation}
\tau_{\rm break}^{-1} \sim \rho_0\,\sigma\,\bar v
\quad\text{with}\quad
\bar v = \sqrt{\frac{k_BT(m_L+m_H)}{m_L m_H}},
\end{equation}
matching the estimate of Ref.~\cite{PaperIII}
($\tau_{\rm break}\approx 0.05$ at our parameters), and giving
the natural scale of the collision integral $\mathcal{C}$.

The collision integral $\mathcal{C}_\alpha$ of
Eq.~\eqref{eq:Calpha} satisfies an H-theorem: the contribution
to the time derivative of the kinetic entropy
$S = -\sum_\alpha\int\!ds\,dv\,\rho_\alpha\log\rho_\alpha$
from $\mathcal{C}_\alpha$ is non-negative, vanishing only at
local Maxwell--Boltzmann equilibrium. The streaming term
$\partial_s(v^{\rm eff}_\alpha\rho_\alpha)$ conserves $S$ on a
periodic domain. The full GHD--Boltzmann dynamics is therefore
monotonically irreversible at the coarse-grained level. This is
qualitatively consistent with the monotonic decay of coherent
kinetic energy $E_{\rm coh}(t)$ reported in Ref.~\cite{PaperIII}
(Sec.~VII therein); a quantitative correspondence between the
kinetic entropy $S(t)$ and $E_{\rm coh}(t)$ would require
computing $S(t)$ directly from the microscopic velocity
distribution, which we have not done.

\subsection{Regime of validity}\label{sec:validity}

The framework requires $\tau_{\rm break}^{-1} \gg k\,c_s$
(strong-coupling hydrodynamic regime). For our parameters,
$k=2\pi/L\approx 0.5$, $c_s\sim 1$ giving
$\tau_{\rm break}\,k\,c_s\approx 0.025$, satisfied
comfortably. The framework \emph{fails} for $m_{\max}\to 1$
where $\tau_{\rm break}\to\infty$ and the integrable
single-species GHD must be used instead; the crossover regime
is parametrically narrow at the densities we consider but
deserves separate study.

\section{Linear regime: sound waves}\label{sec:linear}

We linearise Eq.~\eqref{eq:ghdboltz} about the global
two-species Maxwell--Boltzmann equilibrium at common
temperature $T$ (the equipartition observed in
Ref.~\cite{PaperIII}, Sec.~VII therein, is recovered exactly
by the elastic collision integral $\mathcal{C}_\alpha$),
\begin{equation}
\rho_\alpha^{(0)}(v) = \frac{n_\alpha}{Z_\alpha(T)}
  \exp\!\left[-\frac{m_\alpha v^2}{2 k_B T}\right],\quad
Z_\alpha(T) = \sqrt{\frac{2\pi k_BT}{m_\alpha}},
\label{eq:rhoeq}
\end{equation}
with $n_\alpha = N_\alpha/L = N/(2L)$ for the half-and-half
binary mixture. Writing
$\rho_\alpha(s,v,t) = \rho_\alpha^{(0)}(v) + \delta\rho_\alpha(s,v,t)$
and expanding to first order in $\delta\rho_\alpha$ gives the
linearised GHD--Boltzmann equation
\begin{equation}
\partial_t \delta\rho_\alpha
+ \partial_s\!\left[\tfrac{v}{1-\eta}\,\delta\rho_\alpha
                 + \rho_\alpha^{(0)}(v)\,\delta v^{\rm eff}\right]
= \mathcal{L}_\alpha[\delta\rho_L,\delta\rho_H](v),
\label{eq:ghdlin}
\end{equation}
where $\eta = \sigma\rho_0 = \sigma(n_L+n_H)$,
$\delta v^{\rm eff} = (v - \bar u^{(0)})\,\sigma\delta\rho^{\rm tot}/(1-\eta)^2
- \eta\,\delta\bar u/(1-\eta)$, and
$\mathcal{L}_\alpha$ is the linearised collision operator
acting on $\delta\rho$. At equilibrium $\bar u^{(0)}=0$, so the
two terms simplify to $v\sigma\delta\rho^{\rm tot}/(1-\eta)^2
- \eta\delta\bar u/(1-\eta)$.

\subsection{Conservation laws and hydrodynamic variables}

The linearised collision operator $\mathcal{L}_\alpha$ has four
zero modes corresponding to the four conserved fields of the
binary mixture: per-species number densities $\delta n_L,
\delta n_H$, total momentum density $\delta p$, and total
energy density $\delta e$. Defining
\begin{align}
\delta n_\alpha(s,t) &= \int\!dv\,\delta\rho_\alpha,\\
\delta p(s,t) &= \int\!dv\,v\,(m_L\delta\rho_L + m_H\delta\rho_H),\\
\delta e(s,t) &= \int\!dv\,\tfrac12 v^2\,(m_L\delta\rho_L + m_H\delta\rho_H),
\end{align}
the hydrodynamic projection of Eq.~\eqref{eq:ghdlin} onto these
modes gives a closed $4\!\times\!4$ system at leading order in
gradient.

\subsection{Chapman--Enskog reduction to Euler hydrodynamics}

For a wave at wavenumber $k$ in the hydrodynamic regime
$k\,c_s\,\tau_{\rm break}\ll 1$, we apply the standard
Chapman--Enskog procedure: expand $\delta\rho_\alpha = \delta
\rho_\alpha^{(0)} + \mathrm{Kn}\,\delta\rho_\alpha^{(1)} + \dots$
in the Knudsen number $\mathrm{Kn} \equiv
k\,\lambda_{\rm mfp} \sim k\,c_s\,\tau_{\rm break}$.
At leading order, $\delta\rho_\alpha^{(0)}$ takes the local
Maxwell--Boltzmann form parametrised by the four hydrodynamic
fields $(\delta n_L, \delta n_H, \delta\bar u, \delta T)$,
the only zero modes of $\mathcal{L}_\alpha$. The
non-hydrodynamic modes (relative momentum $p_L-p_H$, relative
temperature $T_L-T_H$) decay at rates $\sim\tau_{\rm break}^{-1}$
and are slaved to the hydrodynamic fields. The solvability
condition of the next-order equation gives the closed system
of macroscopic equations below. The
zeroth-order projection of Eq.~\eqref{eq:ghdlin} onto these
fields gives the linearised two-component Euler equations,
\begin{align}
\partial_t \delta n_\alpha + n_\alpha^{(0)}\partial_s\delta\bar u &= 0,
\quad \alpha\in\{L,H\},\\
\bar m\,\rho_0\,\partial_t \delta\bar u + \partial_s\delta P &= 0,\\
\partial_t \delta T + (\gamma-1)\,T\,\partial_s\delta\bar u &= 0,
\end{align}
with the equation of state given by the Tonks pressure
$P=\rho^{\rm tot} T/(1-\sigma\rho^{\rm tot})$.

\subsection{Sound speed} The adiabatic
sound speed follows from
$c_s^2 = (\partial P/\partial \rho^{\rm tot})_S$ evaluated at
equilibrium, with the entropy held fixed at constant
concentration. Using the 1D ideal-gas adiabatic index
$\gamma=3$ and the average mass per particle
$\bar m = (m_L+m_H)/2$,
\begin{equation}
\boxed{\;c_s^2 = \frac{\gamma\,k_B T}{\bar m\,(1-\eta)^2}
         = \frac{3\,k_B T}{\bar m\,(1-\eta)^2}\;}
\label{eq:cs}
\end{equation}
This is the principal analytical prediction of this section.

For the parameters of Refs.~\cite{PaperII,PaperIII}
($k_BT=1$, $\eta=0.1$, $\bar m=3$ at $m_{\max}=5$):
\begin{equation}
c_s = \sqrt{\frac{3}{3\cdot 0.81}} \approx 1.111.
\label{eq:csnum}
\end{equation}
For $m_{\max}=10$, $\bar m=5.5$ and $c_s\approx 0.821$. For
$m_{\max}=100$, $\bar m=50.5$ and $c_s\approx 0.272$.

\subsection{Comparison with molecular dynamics}

The principal quantitative test is whether
Eq.~\eqref{eq:cs} reproduces the underdamped sound-wave
frequencies of the binary mixture. Ref.~\cite{PaperII} reported
underdamped sound at $m_{\max}=5$; here we extend the
comparison to $m_{\max}\in\{5,10,100\}$ via a dedicated
linear-regime sweep performed in this work.
Table~\ref{tab:cs} compares the GHD--Boltzmann prediction
against the empirical value extracted from this sweep, fitting
the linear-regime
density autocorrelations of Ref.~\cite{PaperII}.

The sweep used $A_\rho=0.10$ (density-perturbation IC), $N=400$,
$\eta=0.10$, $4000$ seeds per $m_{\max}$, $T_{\max}=30$. Fitting
$\mathrm{Re}\,\hat\rho_{k_1}(t) = A e^{-\gamma t}\cos(\omega t
+\varphi)$ extracts $c_s^{\rm MD}=\omega/k$ with $k=2\pi/L$
(Table~\ref{tab:cs}).

\begin{table}[t]
\centering
\caption{Sound speed: GHD--Boltzmann prediction
Eq.~\eqref{eq:cs} versus a dedicated event-driven MD sweep
(this work) with linear-regime density-perturbation IC,
extending the $m_{\max}=5$ result of Ref.~\cite{PaperII}. The
mean error across the three $m_{\max}$ values is $2.5\%$, with
no systematic bias.}
\label{tab:cs}
\begin{tabular}{ccccc}
\hline\hline
$m_{\max}$ & $\bar m$ & $c_s^{\rm GHD}$
                       & $c_s^{\rm MD}$
                       & $c_s^{\rm MD}/c_s^{\rm GHD}$\\
\hline
$5$   & $3.0$  & $1.111$ & $1.151$ & $1.036$\\
$10$  & $5.5$  & $0.821$ & $0.827$ & $1.008$\\
$100$ & $50.5$ & $0.271$ & $0.261$ & $0.965$\\
\hline\hline
\end{tabular}
\end{table}

Across a 20-fold range in mass disparity, the GHD--Boltzmann
sound speed agrees with the microscopic dynamics within
$\sim 4\%$ and with no systematic bias. This is the principal
quantitative validation of the two-species GHD framework
introduced in Sec.~\ref{sec:twospecies}.

\subsection{Damping rate}\label{sec:damping}

The fitted damping rates of Table~\ref{tab:cs} are
$\gamma\approx 0.02$ ($m_{\max}=5,10$) and $\gamma\approx 0.007$
($m_{\max}=100$). These are two orders of magnitude smaller
than the inter-species kinetic-energy exchange rate
$\tau_{\rm break}^{-1}\cdot 4 m_L m_H/(m_L+m_H)^2$ that sets
the non-hydrodynamic spectrum of $\mathcal{L}_\alpha$,
confirming the regime $k\,c_s\,\tau_{\rm break}\ll 1$ assumed in
deriving Eq.~\eqref{eq:cs}. The damping is instead consistent
with viscous Navier--Stokes attenuation,
\begin{equation}
\gamma_{\rm vis} \;\sim\; \tfrac{2}{3}\,\nu_{\rm shear}\,k^2,
\label{eq:gammavis}
\end{equation}
with kinematic viscosity $\nu_{\rm shear}\sim v_{\rm th}
\lambda_{\rm mfp}\approx 0.03$ at our parameters (Sec.~III.D of
Ref.~\cite{PaperIII}), giving $\gamma_{\rm vis}\approx 0.005$ at
$k=2\pi/L=0.5$ --- correct order of magnitude. For a binary mixture the damping
generally has contributions from both shear viscosity
$\nu_{\rm shear}$ and inter-species diffusion $D_{LH}$; the
relative weight depends on the wave structure and on the
species concentration. A first-principles computation of
$\nu_{\rm shear}(m_{\max})$ and $D_{LH}(m_{\max})$ from the
Boltzmann collision integral $\mathcal{C}_\alpha$, and
quantitative comparison with the fitted damping rates, is
deferred to the numerical implementation of
Eq.~\eqref{eq:ghdboltz}.

\section{Density-mode spectrum and $m_{\max}$-universality}
\label{sec:nonlinear}

The principal nonlinear test of the GHD--Boltzmann framework is
the spectral universality reported in
Ref.~\cite{PaperIII}: the density-mode shape
$|\hat\rho_n|/|\hat\rho_2|$ varies by less than $20\%$
between $m_{\max}=5$ and $m_{\max}=100$, despite a 20-fold
variation in the integrability-breaking parameter and despite
the dimensionless step amplitude $U_0/c_s$ varying from $0.45$
(subsonic at $m_{\max}=5$) to $1.85$ (supersonic at $m_{\max}=100$).
A naive estimate of the nonlinear coupling amplitude would
scale as $(U_0/c_s)^2$, predicting an order-of-magnitude
variation; the data show $\sim 20\%$. We propose an explanation
for this saturation.

\subsection{Linear regime: odd-mode universality}

The step initial condition
$u_{\rm coh}(s,0)=U_0\,\mathrm{sgn}[\cos(ks)]$ has continuous
Fourier coefficients
\begin{equation}
\hat u^{(c)}_n(0) = \frac{2U_0}{\pi n}\quad\text{(odd $n$)},
\qquad \hat u^{(c)}_n(0) = 0 \quad\text{(even $n$)}.
\label{eq:uic}
\end{equation}
At first order in $U_0/c_s$, the sound-wave eigenmodes of the
Chapman--Enskog-reduced Euler equations
(Sec.~\ref{sec:linear}) give the coupled system
$\dot{\hat\rho}_n = -ik_n \hat u_n^{(c)}$ and
$\dot{\hat u}_n^{(c)} = -ik_n c_s^2\,\hat\rho_n$. For step IC
with $\hat\rho_n(0)=0$ and $\hat u_n^{(c)}(0)=2U_0/(\pi n)$
for odd $n$, the solution is
$\hat\rho_n(t) = -i\,(\hat u^{(c)}_n(0)/c_s)\,\sin(\omega_n t)\,
e^{-\gamma_n t}$ with $\omega_n = k_n c_s$. The maximum
over time occurs at the first quarter-oscillation
$t_n^{\rm peak} = \pi/(2\omega_n)$ with amplitude
\begin{equation}
\max_t |\hat\rho_n^{(1)}(t)|
= \frac{|\hat u^{(c)}_n|}{c_s}\,e^{-\gamma_n t_n^{\rm peak}}
\;\approx\; \frac{|\hat u^{(c)}_n|}{c_s}
\quad\text{(odd $n$)},
\label{eq:rhomax}
\end{equation}
where the last approximation uses
$\gamma_n t_n^{\rm peak}\ll 1$ for our parameters
($\gamma_n\sim 0.02$, $t_n^{\rm peak}\sim 1$).
The dimensionless ratio is therefore
\begin{equation}
\frac{\max_t|\hat\rho_n^{(1)}|}{\max_t|\hat\rho_1^{(1)}|}
= \frac{|\hat u^{(c)}_n|}{|\hat u^{(c)}_1|} = \frac{1}{n}
\quad\text{(odd $n$)},
\label{eq:linearratio}
\end{equation}
\emph{exactly $m_{\max}$-independent.} The prediction $1/n$ for
odd $n$ is compared directly with the data of
Ref.~\cite{PaperIII} in Table~\ref{tab:nonlin}.

\begin{table}[t]
\centering
\caption{Density-mode ratios versus event-driven MD with step
IC ($U_0=0.5$). For odd $n$ the linear prediction is
$|\hat u^{(c)}_n|/|\hat u^{(c)}_1|=1/n$
(Eq.~\ref{eq:linearratio}), reproduced to $\sim 5\%$. For even
$n$ the linear prediction is exactly zero and the observed
nonzero value is generated by second-order nonlinear coupling
(Sec.~\ref{sec:nonlinear}, subsection on even modes).}
\label{tab:nonlin}
\begin{tabular}{ccccccc}
\hline\hline
& \multicolumn{3}{c}{odd $n$ (linear: $1/n$)}
& \multicolumn{3}{c}{even $n$ (linear: $0$)}\\
$m_{\max}$ & $n{=}3$ & $n{=}5$ & $n{=}7$
           & $n{=}2$ & $n{=}4$ & $n{=}6$ \\
\hline
linear pred. & $0.333$ & $0.200$ & $0.143$
             & $0$     & $0$     & $0$\\
$5$    & $0.333$ & $0.199$ & $0.144$ & $0.260$ & $0.101$ & $0.064$\\
$20$   & $0.330$ & $0.198$ & $0.138$ & $0.312$ & $0.122$ & $0.072$\\
$100$  & $0.330$ & $0.195$ & $0.137$ & $0.278$ & $0.127$ & $0.072$\\
\hline\hline
\end{tabular}
\end{table}

The numerical agreement on the ratio is striking: across a
20-fold range in $m_{\max}$, the odd-mode ratios match the
linear prediction to the third decimal place. The robustness
of the ratio is non-trivial: at $m_{\max}=100$ the dimensionless
amplitude $U_0/c_s = 1.85$ is supersonic, and the absolute
amplitude prediction $\max|\hat\rho_n| = |\hat u^{(c)}_n|/c_s$
breaks down accordingly --- for $n=1$ the predicted
$|\hat u^{(c)}_1|/c_s$ is $0.29$ ($m_{\max}=5$, observed $0.26$,
agreement $91\%$) but $1.17$ ($m_{\max}=100$, observed $0.51$,
agreement $44\%$). What survives the breakdown of the linear
amplitude prediction is the \emph{ratio}, where the
$c_s$-dependent prefactor cancels. This pattern is consistent
with the saturation of nonlinear couplings discussed below.

\subsection{Nonlinear regime: even modes}

The even-$n$ density modes are absent in the linear response
(Eq.~\ref{eq:uic}); they arise from nonlinear coupling of
odd modes via the $\partial_s(\delta\rho\,u)$ term in
continuity. At second order in $U_0/c_s$, a naive perturbative
estimate gives $|\hat\rho_2^{(2)}|/|\hat\rho_1^{(1)}| \sim
U_0/c_s$, which would vary by a factor of four between
$m_{\max}=5$ ($U_0/c_s=0.45$) and $m_{\max}=100$
($U_0/c_s=1.85$). The empirical ratio is essentially constant
across the three values
($0.26, 0.31, 0.28$ for $m_{\max}=5,20,100$).
Notably, the $m_{\max}=100$ entry corresponds to the supersonic
regime $U_0/c_s=1.85$, where the linear-amplitude prediction
for the absolute mode amplitude breaks down by a factor of
$\sim 2$ (Sec.~\ref{sec:nonlinear} above). The persistence of
the spectral-shape ratios across the sub-to-supersonic
transition is itself non-trivial.

We do not derive the even-mode amplitudes from first principles
in this work. At the level of dimensional analysis, the
relevant scales are the coherent-decay time
$\tau_{\rm decay}^{\rm coh}\sim 1$
(Ref.~\cite{PaperIII}, Table~VI),
the linear-peak time $t^*_{\rm peak}(n=2) = 1/(k_2 c_s)$
ranging from $0.9$ ($m_{\max}=5$) to $3.7$ ($m_{\max}=100$),
and the dimensionless step amplitude $U_0/c_s$ ranging from
$0.45$ to $1.85$. The data show $|\hat\rho_2|/|\hat\rho_1|$ in
the range $0.26$--$0.31$, varying by $\sim 20\%$ across the
factor-20 $m_{\max}$ range while the naive perturbative scaling
$\sim U_0/c_s$ would predict a factor-four variation. Producing
this saturation quantitatively from the linearised collision
operator $\mathcal{L}_\alpha$ is left as open work
(Sec.~\ref{sec:discussion}); the principal positive prediction
of this section is the odd-mode ratio
Eq.~\eqref{eq:linearratio}, which the data validate
without free parameters.

\subsection{Consistency with no Burgers shock}

A Burgers shock would have $|\hat\rho_n|/|\hat\rho_1|\to{\rm
const}\sim 1$ at $t\sim t^*_{\rm shock}$. Our prediction
Eq.~\eqref{eq:linearratio} gives $1/n$ in the linear regime,
which decays with $n$, opposite to the Burgers prediction.
Since the nonlinear corrections are suppressed by the
saturation discussed in the preceding subsection,
the linearised GHD--Boltzmann spectrum does not attain the
Burgers shape, and the no-shock result of
Ref.~\cite{PaperIII} is recovered as a consequence of
strong-coupling hydrodynamics with two-species mixing in the
parameter regime tested.

\section{Discussion}\label{sec:discussion}

\subsection{Summary of results}

The paper makes two quantitative predictions without free
parameters and shows that both match molecular-dynamics data:
\emph{(i)} the sound speed
$c_s^2 = 3 k_B T/[\bar m\,(1-\eta)^2]$
(Eq.~\ref{eq:cs}) reproduces the linear-regime oscillation
frequencies of the binary mixture across a 20-fold range in
mass disparity, with error $\lesssim 4\%$ and no systematic
bias (Table~\ref{tab:cs}); \emph{(ii)} the odd-density-mode
ratio $\max|\hat\rho_n|/\max|\hat\rho_1| = 1/n$ for the step
initial condition (Eq.~\ref{eq:linearratio}) is recovered to
within $5\%$ in all configurations tested
(Table~\ref{tab:nonlin}). The first is a standard
hydrodynamic result whose derivation in the present framework
is structurally simple; the second is the explanation for the
empirical density-spectrum universality reported in
Ref.~\cite{PaperIII}.

\subsection{Geometric content and the role of $(M,g)$}

The reduction of GHD on a curved 1D manifold to GHD on a flat
ring of length $L$ (Eq.~\ref{eq:ghdmg}) is a special property of
the one-dimensional setting: in arc-length coordinate the bulk
dynamics is geometrically trivial, and the metric enters only
through the perimeter and the Jacobian relating intrinsic to
arc-length observables. In higher-dimensional analogues the
geometric content of the dynamics is non-trivial \cite{
DoyonSpohnYoshimura2017}, and our reduction does not extend
naively.

\subsection{Quenched versus annealed disorder}
\label{sec:quenched-vs-annealed-disorder}

The mass-disorder framework introduced here treats the disorder
as quenched: each realisation carries a fixed random assignment
of masses, and ensemble averages are performed over disorder
realisations. The collision integral $\mathcal{C}_\alpha$ in
Eq.~\eqref{eq:ghdboltz} captures the self-averaging contribution
that survives this average; the implicit assumption is that
the disorder-averaged dynamics is well-approximated by the
dynamics generated by the average disorder configuration. This
self-averaging holds when the mixing time of the dynamics is
shorter than the time scale over which sample-to-sample
fluctuations build up at a given wavenumber. For our parameters
the inter-species collision rate $\tau_{\rm break}^{-1}\sim 20$
is much faster than the longest sound-wave period, so the
mixing condition is satisfied at the wavenumbers we
investigate. A single-realisation analysis, appropriate for
systems with strong localisation or quenched-disorder-induced
glassiness, requires methods outside the present scope; no such
effects have been reported for binary mass disorder at the
densities and temperatures we consider.

\subsection{Relation to prior work}

The present work sits at the intersection of three nearby
research lines. \emph{(i) Generalised hard rods}: Ferrari and
Olla \cite{HRDiffusive2023} consider hard rods with iid random
lengths drawn from a fixed distribution, deriving Euler-scale
GHD and diffusive corrections. Our setup differs in the
disorder channel (mass, not length) and in admitting a
non-trivial inter-species collision integral; the elastic
binary mass collision Eq.~\eqref{eq:elcoll} has no analogue in
their length-disorder framework.
\emph{(ii) Weak-integrability-breaking GHD}: Durnin, Bhaseen,
and Doyon \cite{DurninDoyon2021}, Bastianello and co-workers
\cite{WeakBreaking2023,BastianelloBoltzmann2020} treat
annealed (stochastic) perturbations of integrable models with
Boltzmann-type collision integrals. Our perturbation is
\emph{quenched} (fixed mass realisation per disorder
configuration); the resulting framework is structurally
similar but the self-averaging assumption
(Sec.~\ref{sec:quenched-vs-annealed-disorder}) is specific to
quenched disorder and requires the mixing condition discussed
there.
\emph{(iii) Mesoscopic GHD with impurities}: Ref.~\cite{
MesoscopicImpurities2023} considers a \emph{single} heavy
impurity in an otherwise integrable gas, deriving its drag
force. Our setup has extensive ($O(N)$) binary disorder,
producing a hydrodynamic regime rather than a single-impurity
limit.

Together with companion papers
\cite{PaperI,PaperII,PaperIII}, the present paper closes a
self-contained four-part program: the equilibrium and
arc-length structure of hard rods on $(M,g)$ (I), the
linear-regime sound-wave restoration under mass disorder (II),
the absence of Burgers shocks in the nonlinear regime (III),
and the GHD-Boltzmann framework that unifies these results into
a single quantitative theory with no free parameters (IV).

\subsection{Open work and follow-up}

Three directions are natural follow-up.
\emph{(i) Damping rate from first principles.} Computing the
shear viscosity $\nu_{\rm shear}(m_{\max})$ as the leading
non-hydrodynamic eigenvalue of $\mathcal{L}_\alpha$ would
upgrade the order-of-magnitude estimate of Sec.~\ref{sec:damping}
to a quantitative prediction comparable with the fitted
damping rates of Table~\ref{tab:cs}.
\emph{(ii) Rigorous treatment of the even-mode saturation.} The
qualitative phenomenology of Sec.~\ref{sec:nonlinear} should
be replaced by a derivation from the linearised collision
operator, yielding the saturated amplitude
$|\hat\rho_2|/|\hat\rho_1|$ and its slow drift across $m_{\max}$
quantitatively.
\emph{(iii) Direct numerical solution.} A numerical
implementation of the two-species GHD-Boltzmann equation
Eq.~\eqref{eq:ghdboltz} on the periodic interval, compared
against the full $\hat\rho_n(t)$ trajectories of
Ref.~\cite{PaperIII}, would test the formalism beyond linear
response and probe the regime where the explicit collision
integral is essential.
\emph{(iv) Microscopic derivation of the dressed velocity.}
Eq.~\eqref{eq:veffalpha} is the central postulate of the
construction (Sec.~\ref{sec:twospecies}) and is currently
justified only by structural analogy with the equal-mass case
and by its empirical success in Sec.~\ref{sec:linear}. A
rigorous derivation from the binary-mixture crossing
dynamics, possibly via a generalisation of the Jepsen mapping
that accounts for the velocity redistribution at unequal-mass
collisions, would close the most significant gap in the
present formalism.

\subsection{Relevance for cold-atom experiments}

Cold-atom realisations of GHD in one-dimensional Bose gases
\cite{Schemmer2019,Malvania2021,Cataldini2022} provide a
natural experimental platform. The principal operational
implication of our results is that a controlled binary mass
disparity --- the cleanest disorder channel in cold-atom
mixture experiments --- modifies the sound speed of the gas in
a quantitatively predictable way (Eq.~\ref{eq:cs}), without
inducing the breakdown of GHD itself. For a $^{6}$Li-$^{7}$Li mixture at $T=100$~nK and packing
$\eta = 0.3$, $\bar m = 6.5\,m_u$ and Eq.~\eqref{eq:cs}
predicts $c_s\approx 7.2$~mm/s. The isotopic mass ratio
$m_H/m_L = 7/6$ is small, however, so the sensitivity
$\partial c_s/\partial m_{\max}$ in Li--Li is only $\sim 1\%$.
A more sensitive test would use a heteronuclear mixture such
as $^{40}$K-$^{87}$Rb or $^{6}$Li-$^{87}$Rb, where the mass
disparity is order $2$-$15$ and Eq.~\eqref{eq:cs} predicts
sound-speed changes of tens of percent, well within reach of
current sound-wave probes in elongated 1D traps
\cite{Cataldini2022}. A follow-up to the present work will
discuss the most accessible parameter regime.

\begin{acknowledgments}
This work was supported by Tecnol\'ogico de Monterrey.
Computations were performed on the xolotl cluster.
\end{acknowledgments}

\begin{thebibliography}{99}

\bibitem{DoyonSpohn2017}
B.~Doyon, H.~Spohn,
J.~Stat.~Mech.\ 073210 (2017).

\bibitem{Doyon2017Lectures}
B.~Doyon,
SciPost Phys.\ Lect.\ Notes \textbf{18} (2020).

\bibitem{BertiniGhdReview2025}
B.~Bertini, A.~De~Luca, A.~Lerose et al.,
Rev.~Mod.~Phys.\ (forthcoming, 2026).

\bibitem{DenardisDoyon2018}
J.~De~Nardis, D.~Bernard, B.~Doyon,
Phys.~Rev.~Lett.\ \textbf{121}, 160603 (2018).

\bibitem{Bulchandani2018}
V.~B.~Bulchandani, R.~Vasseur, C.~Karrasch, J.~E.~Moore,
Phys.~Rev.~B \textbf{97}, 045407 (2018).

\bibitem{GHDperspective2025}
B.~Doyon, S.~Gopalakrishnan, F.~M{\o}ller,
J.~Schmiedmayer, R.~Vasseur,
Phys.~Rev.~X \textbf{15}, 010501 (2025).

\bibitem{BDS1983}
C.~Boldrighini, R.~L.~Dobrushin, Yu.~M.~Sukhov,
J.~Stat.~Phys.\ \textbf{31}, 577 (1983).

\bibitem{HRDiffusive2023}
P.~A.~Ferrari, S.~Olla,
Ann.~Inst.~H.~Poincar\'e (2023) arXiv:2305.13037.

\bibitem{HubnerDoyon2024}
F.~H\"ubner, B.~Doyon,
arXiv:2406.18322 (2024).

\bibitem{DoyonSpohnYoshimura2017}
B.~Doyon, H.~Spohn, T.~Yoshimura,
Nucl.~Phys.~B \textbf{926}, 570 (2018), arXiv:1704.04409.

\bibitem{DurninDoyon2021}
J.~Durnin, M.~J.~Bhaseen, B.~Doyon,
J.~Stat.~Mech.\ (2021) arXiv:2103.11997.

\bibitem{WeakBreaking2023}
A.~Bastianello, A.~De~Luca, B.~Doyon, J.~De~Nardis,
Phys.~Rev.~Res.\ (2023) arXiv:2302.12804.

\bibitem{BastianelloBoltzmann2020}
A.~Bastianello, A.~De~Luca, R.~Vasseur,
J.~Stat.~Mech.\ \textbf{2020}, 124005 (2020).

\bibitem{MesoscopicImpurities2023}
A.~Bastianello, A.~De~Luca, A.~Lerose, T.~Yoshimura,
arXiv:2307.00551 (2023).

\bibitem{Schemmer2019}
M.~Schemmer, I.~Bouchoule, B.~Doyon, J.~Dubail,
Phys.~Rev.~Lett.\ \textbf{122}, 090601 (2019).

\bibitem{Malvania2021}
N.~Malvania, Y.~Zhang, Y.~Le, J.~Dubail, M.~Rigol, D.~S.~Weiss,
Science \textbf{373}, 1129 (2021).

\bibitem{Cataldini2022}
F.~Cataldini, F.~M{\o}ller, M.~Tajik et al.,
Phys.~Rev.~X \textbf{12}, 041032 (2022).

\bibitem{PaperI}
H.~Medel,
in preparation (2026).

\bibitem{PaperII}
H.~Medel,
in preparation (2026).

\bibitem{PaperIII}
H.~Medel,
in preparation (2026).

\end{thebibliography}

\end{document}
